\section{Introduction} \label{sec:DP-OTB_intro}

This chapter focues on the stochastic convex optimization (SCO) problem. As a quick reminder, the problem is formally defined as
\begin{align}
    \min_{\vecx\in \calX} f(\vecx) = \E_{\xi}[F(\vecx,\xi)]
    \label{eqn:DP-OTB:stochopt}
\end{align}
Here, $\vecx$ represents model parameters or weights lying in a convex domain $\calX\subset \RR^d$, $\xi$ is data randomly sampled from an unknown distribution, and $F(\vecx,\xi)$ is the stochastic estimator of loss $f$, which we will assume to be \emph{both convex and smooth} in $\vecx$. Although we do not know the data distribution, we do know the loss $F$, and we have access to an i.i.d. dataset $\calZ_T=(\xi_1,\dots,\xi_T)$ that may have been obtained via user surveys or volunteer tests. Using this information, we would like to solve the optimization problem (\ref{eqn:DP-OTB:stochopt}). The quality of a putative solution $\hat \vecx$ is measured by the \emph{suboptimality gap} defined as $f(\hat \vecx) - f(\vecx_\star)$ for $\vecx_\star \in \argmin_{\vecx\in \calX} f(\vecx)$.

In addition to solving (\ref{eqn:DP-OTB:stochopt}), we also wish to \emph{preserve privacy} for the people who contributed to the dataset $\calZ_T$. To this end, we require our algorithms to be \emph{differentially private} \citep{dwork2006calibrating, dwork2014algorithmic}, which means that replacing any one $\xi_t$ with a different $\xi_t'$ has a negligible effect on the solution $\hat \vecx$. There is a delicate interplay between privacy, dataset size, and solution quality. As $T$ grows, the influence of any individual $\xi_t$ on $\hat \vecx$ decreases, increasing privacy. However, for any finite $T$, one must necessarily leak \emph{some} information about the dataset in order to achieve a nontrivial $f(\hat \vecx)-f(\vecx_\star)$. The goal, then, is to minimize $f(\hat \vecx)-f(\vecx_\star)$ subject to the privacy constraint.

This problem has been well-studied in the literature, and by now the optimal tradeoffs are known and achievable\footnote{We focus on the harder stochastic optimization problem rather than empirical risk minimization, which is also well-studied, e.g. \citep{chaudhuri2011differentially, kifer2012private}.}. In particular, \cite{bassily2019private} exhibits a polytime algorithm that achieves $f(\hat \vecx)-f(\vecx_\star)\le \tilde O(1/\sqrt{T} + \sqrt{d}/\rho T)$ where $\rho$ is a measure of privacy loss (large $\rho$ means less private), and moreover they show that this bound is tight in the worst case. Then, \cite{feldman2020private} provides an improved algorithm that obtains the same guarantee in $O(T)$ time. Further work on this problem considers assumptions on the geometry \citep{asi2021private}, gradient distributions \citep{kamath2021improved}, or smoothness \citep{bassily2020stability, kulkarni2021private}.

All private optimization algorithms we are aware of fall into one of two camps: either they employ some simple pre-processing that ``sanitizes'' the inputs to a non-private optimization algorithm (e.g. the empirically popular DP-SGD of \cite{abadi2016deep}), or they make a careful analysis of the dynamics of their algorithm (e.g. bounding the sensitivity of a single step of stochastic gradient descent). In the former case, the algorithm typically does \emph{not} achieve the optimal convergence rate for stochastic optimization. In the latter case, the algorithm becomes more rigidly tied to the privacy analysis, resulting in delicate ``theory-crafted'' methods that are less popular in practice.

In contrast, in the non-private setting, there is a simple and general technique to produce stochastic optimization algorithms with optimal convergence guarantees: the \emph{online-to-batch conversion} \citep{cesa2004generalization}. This method directly converts any \emph{online convex optimization} (OCO)\citep{shalev2011online, hazan2019introduction, orabona2019modern} algorithm into a stochastic optimization algorithm. OCO is a simple and elegant game-theoretic formulation of the optimization problem that has witnessed an explosion of diverse algorithms and techniques, so that this conversion result immediately implies a vast array of practical optimization algorithms.
In summary: producing \emph{private} stochastic optimization algorithms with optimal convergence rates is currently delicate and difficult, while producing non-private algorithms is essentially trivial.

Our goal is to rectify this issue. To do so, we produce a new \emph{differentially private} online-to-batch conversion. In direct analogy to the non-private conversion, our method converts any OCO algorithm into a \emph{private} stochastic optimization algorithm. After using our conversion, any OCO algorithm that achieves the optimal regret (the standard measure of algorithm quality in OCO), will automatically achieve the optimal suboptimality gap of $\tilde O(1/\sqrt{T} + \sqrt{d}/\rho T)$. Our conversion has additional desirable properties: although convexity is required for the guarantee on suboptimality, it is \emph{not} required for the privacy guarantee, meaning that the method can be easily applied to non-convex tasks (e.g. in deep learning). Further, by largely decoupling the privacy analysis from the algorithm design through this reduction, we can leverage the rich literature of OCO to build private algorithms with new \emph{adaptive} guarantees. For two explicit examples, (1) we construct an algorithm guaranteeing $f(\hat \vecx)-f(\vecx_\star)\le \tilde O(\sigma /\sqrt{T}+ \sqrt{d}/\rho T)$, where $\sigma$ is the apriori unknown standard deviation in the gradient $\nabla F(\vecx,\xi)$, and (2), we construct an algorithm guaranteeing $f(\hat \vecx)-f(\vecx_\star)\le \tilde O(\|\vecx_\star-\vecx_1\|/\sqrt{T} + \sqrt{d}/\rho T) $, where $\vecx_1$ is any user-supplied point. This last guarantee may have applications in private \emph{fine tuning} \citep{li2022private, yu2021differentially, hoory2021learning, kurakin2022toward, mehta2022large}, as the bound automatically improves when the algorithm is provided with a good initialization point.


\section{Preliminaries}

Throughout this chapter, we denote the stochastic loss by $F: \calX\times\Xi \to \RR$, where $\calX\subseteq \RR^d$ is the parameter space (also known as function domain) and $\Xi$ is the data space. In particular, we assume $\calX$ is a general \emph{Banach space} instead of the Euclidean space, meaning that it could be equipped with an arbitrary norm $\|\cdot\|$ beyond the Euclidean norm. We denote its dual norm by $\|\cdot\|_\ast$ and the dual pairing by $\langle \vecx_\ast, \vecx\rangle \coloneqq \vecx_*(\vecx)$. We use $\calZ_T=(\xi_1,\ldots,\xi_T) \in \Xi^T$ to denote an i.i.d. dataset where each $\xi_t$ is sampled i.i.d. from an unknown distribution over $\Xi$.

We will need the following assumptions on the domain $\calX$ and the loss function $F$ for our main results.

\begin{assumption}[Strongly Convex Norm]
\label{as:DP-OTB:sc-norm}
$\|\cdot\|^2$ is $\lambda$-strongly convex w.r.t. $\|\cdot\|$:
\begin{align}
    \|\vecy\|^2 \geq \|\vecx\|^2 + \langle \vecg, \vecy-\vecx\rangle + \frac{\lambda}{2}\|\vecy-\vecx\|^2,
    \qquad \forall \vecx,\vecy\in\RR^d, \vecg \in \partial \|\vecx\|^2.
    \label{eq:DP-OTB:sc-norm}
\end{align}
\end{assumption}

\begin{assumption}[Bounded Domain]
\label{as:DP-OTB:bounded-domain}
Domain $\calX$ of $F$ has diameter at most $D$:
\begin{equation}
    \|\vecx - \vecy\| \le D, 
    \qquad \forall \vecx,\vecy\in\calX.
    \label{eq:DP-OTB:bounded-domain}
\end{equation}
\end{assumption}

\begin{assumption}[Lipschitzness]
\label{as:DP-OTB:Lipschitz}
$F(\vecx,\xi)$ is differentiable and $G$-Lipschitz in $\vecx$:
\begin{equation}
    \|\nabla F(\vecx,\xi)\|_\ast \le G,
    \qquad \forall \vecx\in \calX, \xi\in \Xi.
    \label{eq:DP-OTB:lipschitz}
\end{equation}
\end{assumption}

\begin{assumption}[Smoothness]
\label{as:DP-OTB:smooth}
$F(\vecx,\xi)$ is differentiable and $H$-smooth in $\vecx$: 
\begin{equation}
    \|\nabla F(\vecx,\xi) - \nabla F(\vecy,\xi)\|_* \le H\|\vecx - \vecy\|,
    \qquad \forall \vecx,\vecy\in\calX, \xi\in\Xi.
\end{equation}
\end{assumption}

\begin{assumption}[Bounded Gradient Variance]
\label{as:DP-OTB:sigma-G}
$\nabla F(\vecx,\cdot)$ has variance bounded by $\sigma_G^2$:
\begin{equation}
    \E[\|\nabla f(\vecx) - \nabla F(\vecx,\xi)\|_\ast^2] \le \sigma_G^2,
    \qquad \forall \vecx\in\calX, \xi\in\Xi.
\end{equation}
\end{assumption}

\begin{assumption}[Gradient Difference Variance is Smooth]
\label{as:DP-OTB:sigma-H}
\begin{equation}
    \E[\|[\nabla f(\vecx)-\nabla f(\vecy)] - [\nabla F(\vecx,\xi)-\nabla F(\vecy,\xi)]\|_\ast^2] \leq \sigma_H^2\|\vecx-\vecy\|^2,
    \qquad \forall \vecx,\vecy,\xi.
\end{equation}
\end{assumption}

Note that if $F$ is $G$-Lipschitz and $H$-smooth in $x$, then so it $f$, and Assumption \ref{as:DP-OTB:sigma-G} and \ref{as:DP-OTB:sigma-H} hold with $\sigma_G\leq 2G$ and $\sigma_H\leq 2H$. Moreover, notice that Assumption \ref{as:DP-OTB:bounded-domain} - \ref{as:DP-OTB:sigma-H} depend on the norm $\|\cdot\|$. 


\section{Differential Privacy} 

Differential Privacy (DP) is one of the main focuses of this and the next chapter, and this section provides a comprehensive introduction. At a high level, DP captures how much information is leaked through a randomized algorithm by measuring the output difference of that algorithm if a single input data is changed.

The formal definition of DP hinges on the notion of \emph{neighboring datasets}. Two datasets $\calZ_T=(\xi_1,\ldots,\xi_T)$ and $\calZ'=(\xi_1',\ldots,\xi_T')$ in $\Xi^T$ are said to be neighbors if they differ by exactly one data point, i.e., $|\calZ_T-\calZ_T'| \coloneqq |\{t : \xi_t\ne \xi'_t\}| = 1$. With this notion, we can define:\footnote{Unlike the standard privacy literature, which uses $(\epsilon,\delta)$ for privacy parameters (i.e., $(\epsilon,\delta)$-DP), we use $(\rho,\gamma)$ instead, reserving $\epsilon,\delta$ for convergence metrics.}

\begin{definition}[$(\rho,\gamma)$-DP \citep{dwork2014algorithmic}] A randomized algorithm $\calM: \Xi^T \to \RR^d$ is $(\rho,\gamma)$-differentially private for $\rho,\gamma\ge 0$ if for any neighboring $\calZ_T,\calZ_T'\in \Xi^T$,
\begin{align}
    \P\{\calM(\calZ_T) \in S\} &\le \exp(\rho)\P\{\calM(\calZ_T')\in S\} + \gamma,
    \qquad \forall S\subseteq \RR^d.
    \label{eq:DP-OTB:DP}
\end{align}
\end{definition}

An alternative (and almost equivalent) definition is \emph{R\'enyi differential privacy} (RDP), which allows to compose privacy mechanisms more easily and achieve tighter privacy bounds in certain cases. 

\begin{definition}[$(\alpha,\epsilon)$-RDP \citep{mironov2017renyi}] A randomized algorithm $\calM: \Xi^T \to \RR^d$ is said to be $(\alpha,\rho)$-R\'enyi differentially private for $\alpha>1,\epsilon\ge0$ if for any neighboring datasets $\calZ_T,\calZ_T' \in \Xi^T$,
\begin{align}
    D_{\alpha}(\calM(\calZ_T)\|\calM(\calZ_T'))\le \epsilon,
    \label{eq:DP-OTB:RDP}
\end{align}
where $D_{\alpha}(P\|Q) \coloneqq \frac{1}{\alpha-1}\log \left(\E_{x \sim Q} \left[\frac{p(x)}{q(x)}\right]\right)$ is the \emph{R\'enyi divergence} between two distributions $P,Q$ with corresponding PDFs $p,q$.
\end{definition}

RDP can be easily converted to the usual $(\rho,\gamma)$-DP as follows \citep{mironov2017renyi}: if $\calM$ is $(\alpha, \alpha\rho^2/2)$-RDP for all $\alpha>1$, then it is also $(2\rho\sqrt{\log(1/\gamma)}, \gamma)$-DP for all $\gamma\ge \exp(-\rho^2)$. For this reason, we stick to the notion of RDP throughout this thesis.

Finally, we introduce the definition of \emph{sensitivity}. Almost all methods for ensuring differential privacy involve injecting noise whose scale increases with the sensitivity of the output. In other words, small sensitivity implies high privacy.

\begin{definition}[Sensitivity] 
The sensitivity of a randomized algorithm $\calM: \Xi^T \to (\calX,\|\cdot\|)$, where $(\calX,\|\cdot\|)$ is a Banach space, is defined as
\begin{align}
    \sens(\calM) \coloneqq \sup_{|\calZ_T - \calZ_T'|=1} \|\calM(\calZ_T) - \calM(\calZ_T')\|.
    \label{eq:DP-OTB:sensitivity}
\end{align}
\end{definition}


\subsection{The Tree Mechanism}

At the core of our proposed conversion algorithms, it incorporates a variant of the tree mechanism, which is an advanced privacy mechanism to aggregate the streaming sum of a sequence of outputs. This mechanism is presented in Algorithm \ref{alg:DP-OTB:tree}, whose privacy guarantee is stated in Theorem \ref{thm:DP-O2NC:tree} and proved in Appendix \ref{app:DP-O2NC:tree}. 



Intuitively, the tree mechanism says if each of the algorithms $\calM_1,\ldots,\calM_n$ requires $\sigma^2$ noise for privacy, then the sum $\sum_{i=1}^n\calM_i$ only needs $O(\ln(n)\sigma^2)$ noise for the same level of privacy. With Theorem \ref{thm:DP-O2NC:tree}, we determine the minimal noise required to ensure privacy for Algorithm \ref{alg:DP-O2NC:PONC} in Corollary \ref{cor:DP-O2NC:privacy}, and we evaluate the total noise added by the tree mechanism in Corollary \ref{cor:DP-O2NC:tree-noise}.

\begin{algorithm}[!t]
    \begin{algorithmic}[1]
        \Statex \textbf{Input:} Index $t\in[T]$, global variables $\sigma_1,\ldots,\sigma_T$.
        % \If {$r = 1$}
        %     \State Reset $\textsc{Noise} = \{\}$.
        % \EndIf
        \State If $t=1$, set $\textsc{Noise}\gets\{\}$.
        \State Sample $\xi_t \sim N(0,\sigma_t^2 I)$ and store $\xi_t$ in $\textsc{Noise}$.
        % $\textsc{Noise} \gets \textsc{Noise}\cup\{\xi_t\}$.
        \State Return $\sum_{(\cdot,i)\in \textsc{Node}(t)} \xi_i$.
    
        \vspace{1em}
    
        \Function{Node}{t}
            % \State Set $k\gets 0$, $S\gets$ binary expression of $r$, $L\gets |S|=\lfloor \log_2(r)\rfloor + 1$, and $\textsc{Node}=\{\}$.
            % \For {$i=1,\ldots,L$ such that $S[i]=1$}
            %     \State Set $k' \gets k+2^{L-i}$, store $\textsc{Node}\gets \textsc{Node}\cup\{(k+1,k')\}$, and update $k\gets k'$.
            % \EndFor
            % \State Return $\textsc{Node}$.
            \State \textbf{Initialize:} Set $k\gets 0$ and $S\gets\{\}$.
            \For {$i=0,\ldots,\lceil\log_2(t)\rceil$ while $k<t$}
                \State Set $k'\gets k+2^{\lceil\log_2(t)\rceil-i}$. If $k' \le t$, store $S\gets S\cup\{(k+1),k'\}$ and update $k\gets k'$.
            \EndFor
            \State Return $S$.
        \EndFunction
    \end{algorithmic}
    \caption{Tree Mechanism}
    \label{alg:DP-OTB:tree}
\end{algorithm}

\begin{remark}
\label{rmk:DP-OTB:tree}
To better understand the \textsc{Node} function in Algorithm \ref{alg:DP-OTB:tree}, consider a binary tree of depth $\lceil\log_2(t)\rceil$. $\textsc{Node}(t)$ returns the largest node $n_i$ in each layer $i$, if exists, such that $n_j<n_i\le t$ starting from the top level. As an example, $\textsc{Node}(7)=\{(1,4),(5,6),(7,7)\}$ and $\textsc{Node}(8)=\{(1,8)\}$. In particular, note that $|\textsc{node}(t)| \le \lceil\log_2(t)\rceil \le 2\ln(t)$.
% Note that $\sum_{(i,j)\in\textsc{Node}(t)}\sum_{k=i}^j = \sum_{k=1}^t$.
\end{remark}

Next, we present the main privacy theorem for the tree mechanism as presented in Algorithm \ref{alg:DP-OTB:tree}. 
% Let $\calX$ be a state space and $\calZ^{(1)},\ldots,\calZ^{(n)}$ be dataset spaces. In our case, $\calX$ is the collection of weights $\vecx$ for $f(\vecx,z)$, $\calZ^{(1)}$ is the collection of all possible batches $Z_1^k$ of size $B_1$ defined in Algorithm \ref{alg:DP-O2NC:PONC} and $\calZ^{(i)}$ for $i\ge 2$ is the collection of all possible batches $Z_i^k$ of size $B_2$. 
% We would like to highlight that Theorem \ref{thm:DP-O2NC:tree} holds for any arbitrary $\calX, \calZ, \calM$.

\begin{theorem}
\label{thm:DP-OTB:tree}
Let $\calX$ be state space and $\calZ^{(1)},\ldots,\calZ^{(n)}$ be dataset spaces, and denote $\calZ^{(1:i)}=\calZ^{(1)}\times\dotsm\times\calZ^{(i)}$. Let $\calM^{(i)}:\calX^{i-1}\times\calZ^{(i)} \to \calX$ be a sequence of algorithms for $i\in[n]$, and let $\calA:\calZ^{(1:n)} \to \calX^n$ be the algorithm that, given a dataset $Z_{1:n} \in \calZ^{(1:n)}$, sequentially computes $\vecx_i = \sum_{j=1}^i\calM^{(j)}(\vecx_{1:j-1},Z_i)+\textsc{Tree}(i)$ for $i\in[n]$ and then outputs $\vecx_{1:n}$.
Suppose for all $i\in[n]$ and neighboring $Z_{1:n},Z_{1:n}'\in\calZ^{(1:n)}$, $\|\calM^{(i)}(\vecx_{1:i-1},Z_i)-\calM^{(i)}(\vecx_{1:i-1},Z_i')\| \le s_i$ for all auxiliary inputs $\vecx_{1:i-1}\in\calX^{i-1}$.
Then for all $\alpha>1$, $\calA$ is $(\alpha, \alpha\rho^2/2)$-RDP where
\begin{equation*}
    \rho \le \sqrt{2\ln n} \cdot \max_{b\in[n], i\le b} \frac{s_i}{\sigma_b}.
\end{equation*}
% $\delta\in(0,1)$, $\calA$ is $(\epsilon, \delta)$-DP where
% \begin{equation*}
%     \epsilon \le \sqrt{\log_2(2n)\log 1/\delta} \sup_{b\in[n],i\le b} \frac{\gamma_i}{\sigma_b}.
% \end{equation*}
\end{theorem}

In our case, $\calX$ is the collection of weights $\vecx$ for $f(\vecx,z)$, $\calZ^{(1)}$ is the collection of all possible batches $Z_1^k$ of size $B_1$ defined in Algorithm \ref{alg:DP-O2NC:PONC} and $\calZ^{(t)}$ for $t\ge 2$ is the collection of all possible batches $Z_t^k$ of size $B_2$. Moreover, $\calM^{(1)}$ corresponds to $\grad$ that computes $\vecg_1^k$ and $\calM^{(t)}$ corresponds to $\diff$ that computes $\vecd_t^k$. Upon substituting these specific definition into Theorem \ref{thm:DP-O2NC:tree}, we have the following privacy guarantees.

\begin{lemma}
\label{cor:DP-O2NC:privacy}
Suppose $f(\vecx,z)$ is differentiable and $L$-Lipschitz in $\vecx$ and the domain of $\calA$ is bounded by $D={\delta}/{T}$. Let $B_1,B_2$ satisfies $B_1\ge TB_2/2$. Then for any $\alpha>1, \rho>0$, Algorithm \ref{alg:DP-O2NC:PONC} is $(\alpha,\alpha\rho^2/2)$-RDP if we set the noises $\sigma_{1:T}$ in the tree mechanism as
\begin{equation*}
    % \sigma_1 = \frac{\sqrt{2\ln T}2L}{B_1\rho}, \quad 
    \sigma_t = \sigma := \frac{\sqrt{2\ln T}4dL}{B_2T\rho}.
    % , t=2,\ldots,T.
\end{equation*}
\end{lemma}

\begin{corollary}
\label{cor:DP-O2NC:tree-noise}
Following the assumptions and definitions in Corollary \ref{cor:DP-O2NC:privacy}, for all $t\in [T]$,
\begin{equation*}
    \E\|\textsc{Tree}(t)\|^2 \le 2\ln(t) d\sigma^2 \le \left( \frac{8\ln(T) d^{3/2}L}{B_2T\rho} \right)^2.
\end{equation*}
\end{corollary}